\documentclass[11pt]{book}
\usepackage[margin=0.7in]{geometry}
\usepackage{tikz}
\usepackage{xcolor}
\usepackage{graphicx}
\usepackage{fancyhdr}
\usepackage{microtype}
\usepackage[hidelinks]{hyperref}
\definecolor{ink}{HTML}{17202A}
\definecolor{paper}{HTML}{FFF9ED}
\definecolor{accent}{HTML}{D35400}
\pagecolor{paper}
\color{ink}
\pagestyle{fancy}
\fancyhf{}
\fancyhead[LE,RO]{\small\sffamily\thepage}
\fancyhead[RE]{\small\sffamily\nouppercase{\leftmark}}
\fancyhead[LO]{\small\sffamily\nouppercase{\rightmark}}
\renewcommand{\headrulewidth}{0.4pt}
\newcommand{\panel}[2]{%
  \begin{tikzpicture}
    \node[draw=ink, line width=1.2pt, fill=white, minimum width=0.43\textwidth,
      minimum height=2.7cm, align=center, inner sep=8pt] (p) {\sffamily\Large #1};
    \node[anchor=south west, fill=ink, text=white, font=\scriptsize\sffamily,
      inner sep=3pt] at (p.south west) {#2};
  \end{tikzpicture}%
}
\newcommand{\speech}[1]{\begin{tikzpicture}[baseline=(b.base)]
  \node[draw=ink, fill=white, rounded corners=5pt, inner sep=5pt,
    text width=0.30\textwidth, align=center] (b) {\sffamily #1};
  \draw[ink, thick] (b.south west) -- ++(-5pt,-8pt) -- (b.south east);
\end{tikzpicture}}
\begin{document}
\frontmatter
\begin{titlepage}
  \centering
  \vspace*{1.3in}
  {\fontsize{42}{46}\selectfont\sffamily\bfseries STORY\\[4pt]BOARD\par}
  \vspace{12pt}\textcolor{accent}{\rule{0.72\textwidth}{4pt}}\par
  \vspace{22pt}{\Large A comic-book starter for a small, vivid story\par}
  \vfill
  {\large Author Name\par}
  \vspace{6pt}{\sffamily Issue 01 \textbullet{} \the\year\par}
\end{titlepage}
\chapter*{Credits}
This starter is designed for an editable panel-by-panel workflow. Replace the
sample dialogue, captions, and panel labels with your own work. Keep a short
credits page in every issue so collaborators and artists are easy to identify.
\tableofcontents
\mainmatter
\chapter{The Signal}
\section*{Page 1}
\noindent\panel{A quiet observatory under a wide night sky}{PANEL 1}
\hfill\panel{A console wakes with one red light}{PANEL 2}
\par\medskip
\noindent\speech{Did you hear that?}\hfill
\speech{The receiver has been silent for years.}
\section*{Page 2}
\noindent\panel{Mara reaches for the switch}{PANEL 3}
\hfill\panel{A pattern fills the monitor}{PANEL 4}
\par\medskip
\begin{center}\fbox{\parbox{0.82\textwidth}{\centering\sffamily\itshape
CAPTION: The first message was not a sound. It was a map.}}\end{center}
\chapter{Story Bible}
\section{Characters}
\begin{description}
  \item[Mara Venn] Systems engineer; patient, observant, and too curious to
  leave an unexplained signal alone.
  \item[The Listener] An unknown presence whose messages arrive as geometry.
\end{description}
\section{Places and continuity}
Record recurring locations, props, and visual motifs here. A short story bible
keeps a long-running comic consistent when pages are written out of order.
\chapter*{Issue notes}
\addcontentsline{toc}{chapter}{Issue notes}
\begin{itemize}
  \item Trim size and bleed: choose them before drawing final artwork.
  \item Keep dialogue short enough to read comfortably at the printed size.
  \item Replace the panel placeholders with your own licensed artwork.
\end{itemize}
\backmatter
\chapter*{Colophon}
Typeset with \LaTeX{}. Replace this text with your production credits,
letterer, editor, and edition information.
\end{document}
